\section{The bootstrapping step}

Let $T,T'\subseteq M\otimes F$ be mutually $0.005$-close, let $D\ge2$, and let
$A_1,\dots,A_{D^2}$ be operators on $M$ such that for every $\delta<1$ and
every $\delta$-viable $S\subseteq M$ for $T'$ the span of $A_1S,\dots,A_{D^2}S$
is $\Delta/(1-\delta)^2$-viable for $T'$, where $D^{12}\Delta\le10^{-5}$. These
are the data of Theorem~t:ais of~\cite[lines 1010--1024]{Arad2016RigorousRG}.
For a dimension $q$ write $s=\lfloor q/(2D^2)\rfloor$ and
\[
    \eta(q)=\Bigl(1+\frac{2}{\sqrt{0.96\,s/(8q)}}\Bigr)^{2\dim T'}q\,e^{-s/16}
\]
for the failure probability of random sampling.

\begin{theorem}[One bootstrapping step]
    \label{thm:viable_bootstrap_step}
    \lean{Submodule.exists_isViable_finrank_le_half}
    \leanok
    \uses{thm:viable_random_sampling, thm:subspace_close_trans}
    If $W\subseteq M$ is $0.015$-viable for $T$, has dimension $q\ge4D^2$, and
    $\eta(q)<1$, then some $0.015$-viable subspace for $T$ has dimension at
    most $q/2$.
\end{theorem}

\begin{proof}
    \uses{thm:finite_image_span_dimension}
    \leanok
    By robustness, $W$ is $0.04$-viable for $T'$. Since $\eta(q)<1$, some
    $s$-dimensional subspace $W'\subseteq W$ is $(1-\delta')$-viable for $T'$
    with $\delta'=0.96\,s/(8q)$. The span $S$ of the $A_iW'$ has dimension at
    most $D^2s\le q/2$ and is $\Delta/\delta'^2$-viable for $T'$. Since
    $s\ge q/(4D^2)$, we have $\delta'\ge0.03/D^2$, and $D^{12}\Delta\le10^{-5}$
    with $D\ge2$ gives $\Delta/\delta'^2\le0.0025$. Robustness again shows
    that $S$ is $2(0.0025+0.005)=0.015$-viable for $T$.
\end{proof}

\begin{theorem}[Bootstrapping]
    \label{thm:viable_bootstrap}
    \lean{Submodule.exists_isViable_finrank_lt_of_bootstrap}
    \leanok
    \uses{thm:viable_bootstrap_step}
    If $q_1\ge4D^2$ and $\eta(q)<1$ for every $q\ge q_1$, then some
    $0.015$-viable subspace for $T$ has dimension less than $q_1$.
    This is the abstract argument used in the proof of
    Proposition~prop:albootstrap-dg of~\cite[lines
        1062--1075]{Arad2016RigorousRG}, once the improving family and the
    sampling bound have been supplied. The spatial scope of the physical
    application and the corrected constants are recorded
    in~\cite{gap:alvv16_bootstrap_constants}.
\end{theorem}

\begin{proof}
    \leanok
    The whole factor $M$ is $0.015$-viable. Starting from it, apply the
    previous theorem as long as the dimension is at least $q_1$; each step
    halves the dimension.
\end{proof}

\begin{theorem}[A quantitative bootstrap bound]
    \label{thm:viable_bootstrap_quantitative_dimension}
    \lean{Submodule.exists_isViable_finrank_lt_rank_mul_pow}
    \leanok
    \uses{def:subspace_close, def:subspace_viable}
    Let $M,F$ be finite-dimensional complex inner-product spaces, and let
    $T,T'\subseteq M\otimes F$ be mutually $0.005$-close. Put
    $r=\dim T'$, and assume $r\ge1$. Let $D\ge2$ be an integer and
    $\Delta\in\R$ satisfy $D^{12}\Delta\le10^{-5}$. Suppose a fixed family
    $A_1,\ldots,A_{D^2}$ of linear operators on $M$ has the following property:
    for every real $\delta<1$ and every $\delta$-viable subspace
    $S\subseteq M$ for $T'$, the space $\operatorname{span}_i A_iS$ is
    $\Delta/(1-\delta)^2$-viable for $T'$.
    Then there is a $0.015$-viable subspace $W\subseteq M$ for $T$ with
    $\dim W<rD^{24}$.

    This is a conditional form of the quantitative argument
    in~\cite[lines 1062--1075]{Arad2016RigorousRG}. Applying it to a spin
    chain requires an improving family for the chosen interval, whose role
    is explained in~\cite{gap:alvv16_bootstrap_constants}.
\end{theorem}
\begin{proof}\leanok
    \uses{thm:viable_bootstrap}
    Let $q\ge rD^{24}$ and $s=\lfloor q/(2D^2)\rfloor$.
    Since $r\ge1$ and $D\ge2$, we have $q\ge4D^2$. The integer quotient
    and the covering-net factor satisfy
    \begin{align}
        s                               & \ge\frac{q}{4D^2}, &
        1+\frac{2}{\sqrt{0.96\,s/(8q)}} & \le21D.
        \label{eq:bootstrap_quantitative_sample}
    \end{align}
    The bound $q\ge rD^{24}$ also gives
    \begin{align}
        \bigl(2r\log(21D)+\log q\bigr)256D^2\le2q.
        \label{eq:bootstrap_quantitative_log_budget}
    \end{align}
    Substituting~(\ref{eq:bootstrap_quantitative_sample})
    and~(\ref{eq:bootstrap_quantitative_log_budget}) into the sampling
    failure probability yields
    \begin{align}
        \eta(q) & \le\exp\bigl(2r\log(21D)+\log q-s/16\bigr)
        \le\exp\bigl(-q/(128D^2)\bigr)<1.
        \label{eq:bootstrap_quantitative_failure}
    \end{align}
    Apply Theorem~\ref{thm:viable_bootstrap} with threshold $rD^{24}$.
\end{proof}

\begin{theorem}[Enlargement at a prescribed accuracy]
    \label{thm:viable_enlargement_inverse_scale}
    \lean{Submodule.viable_enlargement_of_close_improving_target}
    \leanok
    \uses{def:subspace_viable, def:subspace_close}
    Let $\mathcal H,\mathcal M,\mathcal E$ be finite-dimensional complex
    inner-product spaces, and fix an isometry
    $u:\mathcal H\to\mathcal M\otimes\mathcal E$ onto the tensor product.
    Let $G,T\subseteq\mathcal H$, let $1\le m\le D$ and $b\ge0$ be integers,
    and suppose that $G,T$ are mutually $1/(200m)$-close.
    Let $W_0\subseteq\mathcal M$ be $0.015$-viable for $uG$, with $\dim W_0\le b$.
    Suppose that $A_0,\ldots,A_{D^2-1}$ act on $\mathcal M$ and that
    $D^{12}\Delta\le10^{-5}$. Assume that, for every real $\delta<1$ and
    every $\delta$-viable subspace $S$ for $uT$, the space
    $\sum_j A_jS$ is $\Delta/(1-\delta)^2$-viable for $uT$.
    Then $V=\sum_j A_jW_0$ satisfies
    \begin{align}
        \dim V&\le bD^2, & V&\text{ is }\frac1{2m}\text{-viable for }uG.
        \label{eq:viable_inverse_scale}
    \end{align}
    This is a quantitative consequence of the closeness transfers in
    \cite[Lemma~l:close and proof of Theorem~4.9]{Arad2016RigorousRG}.
\end{theorem}
\begin{proof}
    \leanok
    \uses{thm:subspace_close_trans, thm:viable_coordinates, thm:finite_image_span_dimension}
    Since $1/(200m)\le0.005$, the first closeness transfer makes $W_0$
    $0.04$-viable for $uT$. The improving estimate and $m\le D^{12}$ give
    \begin{align*}
        \frac{\Delta}{(1-0.04)^2}&\le\frac1{8m}, &
        2\left(\frac1{8m}+\frac1{200m}\right)&=\frac{13}{50m}\le\frac1{2m}.
    \end{align*}
    For $\Delta\le0$ the first inequality is immediate; for $\Delta>0$
    it follows from $m\Delta\le10^{-5}$.
    Transfer back to $uG$. Finally,
    $\dim\sum_j A_jW_0\le\sum_j\dim A_jW_0\le D^2\dim W_0$.
\end{proof}

\begin{lemma}[The ceiling scale]
    \label{lem:ceil_quadratic_scale}
    \lean{Real.ceil_quadratic_scale_bounds}
    \leanok
    If $A,v\ge1$ and $0<g\le1$, then
    \begin{align*}
      Av^2/g\le\lceil Av^2/g\rceil
      \le(A+1)v^2/g,
      \qquad A\le\lceil Av^2/g\rceil.
    \end{align*}
\end{lemma}
\begin{proof}
    \leanok
    The ceiling lies between $Av^2/g$ and $Av^2/g+1$.
    Use $v^2/g\ge1$ to absorb the additive term.
\end{proof}

\begin{lemma}[One dyadic change of ceiling]
    \label{lem:ceil_quadratic_scale_half_step}
    \lean{Real.ceil_quadratic_scale_half_step}
    \leanok
    Suppose $A,u\ge100$, $v\ge1$, $v\le u+1$, and
    $0<g,h\le1$ with $g\le2h$. Put
    $\ell=\lceil Au^2/g\rceil$ and $\ell'=\lceil Av^2/h\rceil$.
    Then $A\le\ell,\ell'$ and $\ell'\le(21/10)\ell$.
\end{lemma}
\begin{proof}
    \leanok
    \uses{lem:ceil_quadratic_scale}
    Use $A+1\le(101/100)A$, $v\le(101/100)u$ and $1/h\le2/g$.
    The resulting factor $2(101/100)^3$ is smaller than $21/10$.
\end{proof}

\begin{theorem}[One dyadic change of count]
    \label{thm:log_polynomial_filter_count_half_step}
    \lean{Real.log_polynomial_filter_count_half_step}
    \leanok
    Fix an integer $m$ with $\log m\ge100$ and a real $A\ge100$.
    For $0<g\le1$, define
    \begin{align*}
      u_g&=1+\log m+\log(1/g),&
      \ell_g&=\lceil Au_g^2/g\rceil,&
      D_g&=(m\ell_g^2)^{22\ell_g}.
    \end{align*}
    If $0<h\le1$ and $g\le2h$, then
    \begin{align*}
      100\le\log D_g,
      \qquad\log D_h\le\frac{13}{6}\log D_g.
    \end{align*}
\end{theorem}
\begin{proof}
    \leanok
    \uses{lem:ceil_quadratic_scale_half_step}
    Since $\log2\le1$, one has $u_h\le u_g+1$ and therefore
    $\ell_h\le(21/10)\ell_g$. Set $B_g=\log m+2\log\ell_g$.
    The inequality $\log(21/10)\le11/10$ gives
    $B_h\le B_g+11/5\le(511/500)B_g$, using $B_g\ge100$.
    Finally $\log D_g=22\ell_gB_g$ and
    $(21/10)(511/500)<13/6$.
\end{proof}
